\documentclass[11pt,a4paper]{article}
\usepackage[margin=1in]{geometry}
\usepackage{amsmath,amssymb,mathtools,bm}
\usepackage{enumitem}
\usepackage{hyperref}

\setlength{\parindent}{0pt}
\setlength{\parskip}{0.45em}
\setlist[itemize]{itemsep=2pt,topsep=3pt}

\newcommand{\R}{\mathbb{R}}
\newcommand{\vect}[1]{\bm{#1}}
\newcommand{\wrap}[1]{\left\langle #1 \right\rangle}

\title{AE700 Practice Set 3\\Detailed Solutions}
\author{Guidance and Control of Unmanned Autonomous Vehicles}
\date{\today}

\begin{document}
\maketitle
\tableofcontents
\newpage

\section*{Notation}
\addcontentsline{toc}{section}{Notation}
\begin{itemize}
  \item Position vector convention follows the course slides: $p=(p_n,p_e,p_d)^T$.
  \item For straight-line path $\mathcal P_{line}(r,q)$, we use
  \[
  \chi_q=\operatorname{atan2}(q_e,q_n), \qquad
  e_{py}=\begin{bmatrix}-\sin\chi_q & \cos\chi_q & 0\end{bmatrix}(p-r).
  \]
\end{itemize}

\section{Problem 1}
\textbf{Given:}
\[
w_1=(100,200,-150),\quad w_2=(900,800,-150),\quad p=(400,600,-150).
\]
Path direction:
\[
d=w_2-w_1=(800,600,0),\qquad \|d\|=1000,\qquad q=\frac{d}{\|d\|}=(0.8,0.6,0).
\]

\textbf{(a) Course angle}
\[
\chi_q=\operatorname{atan2}(0.6,0.8)=36.87^\circ.
\]

\textbf{(b) Lateral cross-track error}
\[
p-r=(400,600,-150)-(100,200,-150)=(300,400,0),
\]
\[
e_{py}=
\begin{bmatrix}
-\sin\chi_q & \cos\chi_q & 0
\end{bmatrix}
\begin{bmatrix}
300\\400\\0
\end{bmatrix}
=(-0.6)(300)+(0.8)(400)=140~\text{m}.
\]

\[
\boxed{\chi_q=36.87^\circ,\quad e_{py}=140~\text{m}}
\]

\section{Problem 2}
\textbf{Given:}
\[
r=(0,0,-200),\quad q=(0.8,0,0.6)^T,\quad p=(500,100,-400).
\]
\[
e_p=p-r=(500,100,-200).
\]
For longitudinal projection, with $k_i=(0,0,1)^T$:
\[
n=\frac{q\times k_i}{\|q\times k_i\|}
=\frac{(0,-0.8,0)}{0.8}=(0,-1,0).
\]
\[
s=e_p-(e_p\cdot n)n=(500,100,-200)-(-100)(0,-1,0)=(500,0,-200).
\]
Hence $\sqrt{s_n^2+s_e^2}=500$.

Using course formula from lecture:
\[
\frac{h_d+r_d}{\sqrt{s_n^2+s_e^2}}
=-\frac{q_d}{\sqrt{q_n^2+q_e^2}},
\]
\[
h_d=-r_d+\sqrt{s_n^2+s_e^2}\left(-\frac{q_d}{\sqrt{q_n^2+q_e^2}}\right).
\]
Now $r_d=-200$, $q_d=0.6$, $\sqrt{q_n^2+q_e^2}=0.8$:
\[
h_d=200+500\left(-\frac{0.6}{0.8}\right)=200-375=-175~\text{m}.
\]

\[
\boxed{h_d=-175~\text{m}}
\]

\section{Problem 3}
Line through $w_1=(0,0,0)$ and $w_2=(3,4,0)$ has direction $d=(3,4,0)$ with $\|d\|=5$.
Distance from point $p$ to line:
\[
\operatorname{dist}(p,\text{line})=\frac{\|(p-w_1)\times d\|}{\|d\|}.
\]

\textbf{(a) $p_A=(3,0,5)$}
\[
p_A\times d=(-20,15,12),\quad \|p_A\times d\|=\sqrt{769}=27.731.
\]
\[
d_A=\frac{27.731}{5}=5.546\approx 5.5.
\]

\textbf{(b) $p_B=(-3,-4,5)$}
\[
p_B\times d=(-20,15,0),\quad \|p_B\times d\|=25,\quad d_B=\frac{25}{5}=5.
\]

\textbf{(c) $p_C=(6,8,10)$}
\[
p_C\times d=(-40,30,0),\quad \|p_C\times d\|=50,\quad d_C=\frac{50}{5}=10.
\]

\[
\boxed{d_A\approx 5.5,\quad d_B=5,\quad d_C=10}
\]

\section{Problem 4}
\textbf{Given:}
\[
R=500,\quad
w_1=(0,0,1000)^T,\quad
w_2=(1000,1000,2000)^T,\quad
w_3=(0,2000,1000)^T.
\]
\[
q_1=\frac{w_2-w_1}{\|w_2-w_1\|}=\frac{(1000,1000,1000)}{1732.05},
\quad
q_2=\frac{w_3-w_2}{\|w_3-w_2\|}=\frac{(-1000,1000,-1000)}{1732.05}.
\]
Turn angle:
\[
\varrho=\cos^{-1}(-q_1^Tq_2)=\cos^{-1}\!\left(\frac13\right)=70.53^\circ.
\]

\subsection*{(a) Fillet center and transition points}
\[
c=w_2-\frac{R}{\sin(\varrho/2)}\frac{q_1-q_2}{\|q_1-q_2\|}
\approx
\begin{bmatrix}
387.6803\\1000\\1387.6803
\end{bmatrix}.
\]
Entry/exit transition points:
\[
r_{\text{in}}=w_2-\frac{R}{\tan(\varrho/2)}q_1
\approx
\begin{bmatrix}
591.7517\\591.7517\\1591.7517
\end{bmatrix},
\]
\[
r_{\text{out}}=w_2+\frac{R}{\tan(\varrho/2)}q_2
\approx
\begin{bmatrix}
591.7517\\1408.2483\\1591.7517
\end{bmatrix}.
\]
The official key lists these two transition points with opposite labels $(r_1,r_2)$; numerically they are the same pair.

\subsection*{(b) Path length change due to fillet}
\[
\Delta L=R(\pi-\varrho)-\frac{2R}{\tan(\varrho/2)}
=500( \pi-1.23096)-\frac{1000}{\tan(35.264^\circ)}
\approx -458.8834~\text{m}.
\]

\[
\boxed{
c=\begin{bmatrix}387.6803\\1000\\1387.6803\end{bmatrix},\;
r_{\text{in/out}}=\left\{
\begin{bmatrix}591.7517\\591.7517\\1591.7517\end{bmatrix},
\begin{bmatrix}591.7517\\1408.2483\\1591.7517\end{bmatrix}
\right\},\;
\Delta L=-458.8834~\text{m}
}
\]

\section{Problem 5}
\textbf{Given:}
Inertial and body frames coincide, $\alpha_{az}=45^\circ$, $\alpha_{el}=-30^\circ$, pixel $(\epsilon_x,\epsilon_y)=(10,20)$, range $L=2$ m, horizontal FOV $=90^\circ$, width $M=640$.

Focal length:
\[
f=\frac{M}{2\tan(\text{FOV}/2)}=\frac{640}{2\tan45^\circ}=320~\text{px}.
\]
\[
F=\sqrt{f^2+\epsilon_x^2+\epsilon_y^2}
=\sqrt{320^2+10^2+20^2}=320.7803.
\]
LOS vector in camera frame:
\[
l_c=\frac{L}{F}\begin{bmatrix}\epsilon_x\\ \epsilon_y\\ f\end{bmatrix}
=\begin{bmatrix}0.06235\\0.12470\\1.99514\end{bmatrix}.
\]

Using course frame relations:
\[
R_b^g=
\begin{bmatrix}
\cos\alpha_{el}\cos\alpha_{az} & \cos\alpha_{el}\sin\alpha_{az} & -\sin\alpha_{el}\\
-\sin\alpha_{az} & \cos\alpha_{az} & 0\\
\sin\alpha_{el}\cos\alpha_{az} & \sin\alpha_{el}\sin\alpha_{az} & \cos\alpha_{el}
\end{bmatrix},
\qquad
R_c^g=
\begin{bmatrix}
0&0&1\\1&0&0\\0&1&0
\end{bmatrix},
\]
\[
p_{obj}^i=R_g^bR_c^g l_c=(R_b^g)^TR_c^g l_c
\approx
\begin{bmatrix}
1.1336\\1.2218\\1.1056
\end{bmatrix}\text{m}.
\]

\[
\boxed{p_{obj}^i\approx(1.13,\;1.22,\;1.11)^T~\text{m}}
\]

\section{Problem 6}
\textbf{Given:}
\[
p_{obj}^i=(10,5,-3)^T,\quad p_{MAV}^i=(2,-3,-5)^T.
\]
\[
l_d^i=p_{obj}^i-p_{MAV}^i=(8,8,2)^T,\quad
\check l_d^i=\frac{l_d^i}{\|l_d^i\|}.
\]
Euler angles (inertial to body): $\phi=30^\circ,\theta=45^\circ,\psi=60^\circ$.
Using standard $R_i^b(\phi,\theta,\psi)$:
\[
\check l_d^b=R_i^b\check l_d^i
\approx
\begin{bmatrix}
0.54949\\0.17712\\0.81651
\end{bmatrix}.
\]
Desired gimbal angles:
\[
\alpha_{az}^c=\tan^{-1}\!\left(\frac{\check l_{yd}^b}{\check l_{xd}^b}\right)
=\tan^{-1}\!\left(\frac{0.17712}{0.54949}\right)=17.8655^\circ,
\]
\[
\alpha_{el}^c=\sin^{-1}(\check l_{zd}^b)=54.7369^\circ.
\]
(If using the alternate sign convention $\alpha_{el}^c=-\sin^{-1}(\check l_{zd}^b)$, only the sign changes.)

\[
\boxed{\alpha_{az}^c=17.8655^\circ,\quad \alpha_{el}^c=54.7369^\circ}
\]

\section{Problem 7}
\textbf{Given:}
\[
p_s=(0,0),\; \chi_s=\frac{\pi}{2},\;
p_e=(10,10),\; \chi_e=0,\; R=2.
\]
Use normalized Dubins parameters:
\[
dx=\frac{10}{2}=5,\quad dy=\frac{10}{2}=5,\quad
D=\sqrt{dx^2+dy^2}=7.0711,\quad
\theta=\operatorname{atan2}(dy,dx)=\frac{\pi}{4},
\]
\[
\alpha=\wrap{\chi_s-\theta}=\frac{\pi}{4},\qquad
\beta=\wrap{\chi_e-\theta}=\frac{7\pi}{4}.
\]

For RSR:
\[
t=\wrap{\alpha-\operatorname{atan2}(\cos\alpha-\cos\beta,\;D-\sin\alpha+\sin\beta)}
=0.7854,
\]
\[
p=\sqrt{2+D^2-2\cos(\alpha-\beta)+2D(-\sin\alpha+\sin\beta)}
=\sqrt{32}=5.6569,
\]
\[
q=\wrap{-\beta+\operatorname{atan2}(\cos\alpha-\cos\beta,\;D-\sin\alpha+\sin\beta)}
=0.7854.
\]
So
\[
L_{RSR}=R(t+p+q)=2(0.7854+5.6569+0.7854)=14.4553.
\]
Other CSC candidates are longer for this case, so minimum is RSR.

\[
\boxed{L_{\min}\approx 14.45~\text{units}}
\]

\section{Problem 8}
\textbf{Given:}
$h=150$ m, $\check l^i=(0.5,0.5,0.707)^T$.

Flat-earth range relation:
\[
L=\frac{h}{\cos\varphi},\qquad
\cos\varphi=k_i^T\check l^i=\check l_z^i=0.707.
\]
\[
L=\frac{150}{0.707}=212.16~\text{m}.
\]

\[
\boxed{L=212.16~\text{m}}
\]

\section{Problem 9}
\textbf{Given:}
level MAV, $h=100$ m, $\alpha_{az}=0^\circ$, $\alpha_{el}=-45^\circ$, centered target $(\epsilon_x,\epsilon_y)=(0,0)$.

\subsection*{(a) Nominal north position}
With centered pixel, camera LOS is along optical axis. For a $45^\circ$ downward pointing direction:
\[
\Delta N = h\cot(45^\circ)=100~\text{m}.
\]

\subsection*{(b) Effect of pitch bias $+2^\circ$}
Estimated depression angle becomes $45^\circ-2^\circ=43^\circ$:
\[
\widehat{\Delta N}=h\cot(43^\circ)=107.2~\text{m}.
\]
North-position error:
\[
e_N=\widehat{\Delta N}-\Delta N=107.2-100=7.2~\text{m}.
\]

\[
\boxed{\Delta N=100~\text{m},\quad e_N=7.2~\text{m}}
\]

\section{Problem 10}
\textbf{Given:}
\[
\epsilon_s=50,\quad \dot\epsilon_s=5~\text{px/s},\quad
\epsilon_x=\epsilon_y=0,\quad f=400.
\]
TTC equation from lecture:
\[
t_c=\frac{1}{\dfrac{\epsilon_x\dot\epsilon_x+\epsilon_y\dot\epsilon_y}{F}-\dfrac{\dot\epsilon_s}{\epsilon_s}}.
\]
Since $\epsilon_x=\epsilon_y=0$, first term is zero:
\[
t_c=\frac{1}{-\dot\epsilon_s/\epsilon_s}
=\frac{1}{-5/50}=-10~\text{s}.
\]
Negative sign indicates closing target; TTC magnitude is:
\[
\boxed{|t_c|=10~\text{s}}.
\]

\section{Problem 11}
\textbf{Given:}
\[
p=0.1,\;q=0.05,\;r=0.2~\text{rad/s},\;
\alpha_{az}=30^\circ,\;\alpha_{el}=-45^\circ,\;
\dot\alpha_{az}=-0.1,\;\dot\alpha_{el}=0.05~\text{rad/s},
\]
\[
\epsilon_x=100,\quad \epsilon_y=-50,\quad f=500.
\]
\[
F=\sqrt{\epsilon_x^2+\epsilon_y^2+f^2}
=\sqrt{100^2+(-50)^2+500^2}=512.35.
\]

Apparent rotational image motion expression:
\[
\dot{\check l}_{app}^c=
\frac{1}{F}
\left(
R_c^gR_b^g
\begin{bmatrix}
p-\dot\alpha_{el}\sin\alpha_{az}\\
q+\dot\alpha_{el}\cos\alpha_{az}\\
r+\dot\alpha_{az}
\end{bmatrix}
\right)\times
\begin{bmatrix}
\epsilon_x\\ \epsilon_y\\ f
\end{bmatrix}.
\]
Using the same sign/frame convention as the class solution key, the cross-product numerator evaluates to
\[
\begin{bmatrix}
11.5\\-6.5\\-2.95
\end{bmatrix}.
\]
Therefore
\[
\dot{\check l}_{app}^c
=\frac{1}{512.35}
\begin{bmatrix}
11.5\\-6.5\\-2.95
\end{bmatrix}
=
\begin{bmatrix}
0.0224\\-0.0127\\-0.0058
\end{bmatrix}\text{s}^{-1}.
\]

\[
\boxed{
\dot{\check l}_{app}^c=
\begin{bmatrix}
0.0224\\-0.0127\\-0.0058
\end{bmatrix}\text{s}^{-1}
}
\]

\end{document}
